\documentclass[11pt,a4paper]{article}
\usepackage[utf8]{inputenc}
\usepackage[T1]{fontenc}
\usepackage{geometry}
\usepackage{graphicx}
\usepackage{hyperref}
\usepackage{listings}
\usepackage{xcolor}
\usepackage{longtable}
\usepackage{booktabs}
\usepackage{amsmath}
\usepackage{calc}
\usepackage{amsthm}
\usepackage{amssymb}
\usepackage{enumitem}
\usepackage{fancyhdr}
\usepackage{abstract}
\usepackage{tocloft}

% Page setup
\geometry{
    a4paper,
    left=25mm,
    right=25mm,
    top=30mm,
    bottom=30mm
}

% Code listing style
\definecolor{codegray}{rgb}{0.5,0.5,0.5}
\definecolor{codepurple}{rgb}{0.58,0,0.82}
\definecolor{codegreen}{rgb}{0,0.6,0}
\definecolor{backcolor}{rgb}{0.95,0.95,0.92}

\lstdefinestyle{mystyle}{
    backgroundcolor=\color{backcolor},
    commentstyle=\color{codegreen},
    keywordstyle=\color{magenta},
    numberstyle=\tiny\color{codegray},
    stringstyle=\color{codepurple},
    basicstyle=\ttfamily\footnotesize,
    breakatwhitespace=false,
    breaklines=true,
    captionpos=b,
    keepspaces=true,
    numbers=left,
    numbersep=5pt,
    showspaces=false,
    showstringspaces=false,
    showtabs=false,
    tabsize=2,
    frame=single
}

\lstset{style=mystyle}

% Header and footer
\pagestyle{fancy}
\fancyhf{}
\rhead{XCS Protocol v1.0}
\lhead{Draft Specification}
\rfoot{Page \thepage}

% Hyperref setup
\hypersetup{
    colorlinks=true,
    linkcolor=blue,
    filecolor=magenta,
    urlcolor=cyan,
    pdftitle={XCS: XRPL Credential Service Protocol},
    pdfauthor={XCS Protocol Team},
    pdfsubject={Blockchain Credentials},
    pdfkeywords={XRPL, Credentials, Blockchain, Protocol}
}

% Title and author information
\title{\Large \textbf{XCS: XRPL Credential Service Protocol}}
\author{Thomas Hussenet thomas@xrpl-commons.org}
\date{Draft Specification v1.0 \\ \today}

\begin{document}

\maketitle

\begin{abstract}
\noindent
XCS (XRPL Credential Service) is an open, neutral, and interoperable credentials registry protocol built on top of the XRP Ledger's native Credentials amendment. By adopting a minimalist design philosophy—thin protocol, rich applications—XCS provides a robust foundation for verifiable credentials while maintaining maximum flexibility for implementers. The protocol establishes core primitives for schema registration, credential lifecycle management, and composable data structures, while deliberately delegating privacy, storage, and advanced features to the application layer. This paper presents the specification, rationale for the design, and implementation considerations of the protocol for XCS.
\end{abstract}

\tableofcontents
\newpage

\section{Introduction}

\subsection{Background}

The digital transformation of identity and credentials faces fundamental challenges: centralization creates single points of failure, proprietary systems lack interoperability, and privacy concerns arise from unnecessary data exposure. While blockchain technology offers solutions through decentralization and cryptographic verification, existing implementations often suffer from over-engineering, attempting to solve all problems at the protocol layer.

The XRP Ledger, with its recent Credentials amendment[1], provides native support for credential issuance, acceptance, and revocation. This creates an opportunity to build a lightweight protocol layer that leverages these primitives while maintaining the flexibility needed for diverse use cases.

\subsection{Motivation}

Current credential systems face several limitations:

\begin{itemize}
    \item \textbf{Vendor lock-in}: Credentials issued by one platform cannot easily be verified by another
    \item \textbf{Lack of composability}: Credentials cannot build upon or reference each other
    \item \textbf{Over-specification}: Protocols that try to solve everything become too complex for wide adoption
    \item \textbf{Privacy vs. Transparency}: Difficulty balancing public verifiability with private data
\end{itemize}

XCS addresses these challenges through a deliberately minimal protocol that provides just enough structure for interoperability while pushing complexity to the edges.

\section{Design Philosophy}

\subsection{Core Principles}

\begin{description}[leftmargin=!,labelwidth=\widthof{\bfseries Composability:}]
    \item[Minimalism:] The protocol includes only what must be coordinated on-chain. Every additional feature is scrutinized for whether it truly belongs at the protocol layer.
    
    \item[Composability:] Schemas can inherit and extend from each other, enabling building blocks for complex credentials without protocol complexity.
    
    \item[Neutrality:] The protocol has no opinion on credential contents, verification methods, or use cases. It provides primitives, not policies.
    
    \item[Openness:] Anyone can create schemas or issue credentials. No gatekeepers, no permission required.
\end{description}

\subsection{Layer Separation}

The protocol explicitly separates concerns between what must be coordinated on-chain and what can be implemented independently by applications.

\subsubsection{Protocol Layer (On-chain)}

The protocol layer contains only the essential primitives that require global consensus. Schema registration and discovery provides a universal registry where credential structures are defined and referenced by unique identifiers, enabling any party to verify credential formats independently. Credential lifecycle management—creation, acceptance, and revocation—leverages the XRP Ledger's native Credentials amendment, inheriting the underlying blockchain's security while maintaining simplicity.

The basic type system offers primitive types that make up complex structures, balancing expressiveness with simplicity. Inheritance rules enable schema composition and reuse through a well-defined mechanism for extending existing schemas, including multiple inheritance and conflict resolution.

\subsubsection{Application Layer (Off-chain)}

The application layer handles all functionality that does not require global consensus. Storage implementation remains flexible—applications choose between IPFS for public data, private databases for sensitive information, or hybrid approaches. The protocol specifies only how to reference the data, not how to store it.

Privacy features like zero-knowledge proofs and selective disclosure operate at the application level, allowing implementations to adapt to varying requirements across jurisdictions and use cases. Advanced revocation logic extends simple revocation with suspension, expiration, reason codes, and conditional reactivation based on specific application needs.

W3C compatibility adapters bridge XCS with existing Verifiable Credentials infrastructure when needed, without forcing compliance at the protocol level. Verification APIs support various strategies, from open public verification to authenticated private checks. Metadata and descriptions provide human-readable context, internationalization, and rich user interfaces without protocol modifications.

\vspace{1em}
This separation ensures that the protocol remains simple, auditable, and resistant to feature creep while enabling rich functionality at the application layer. The protocol provides coordination; applications provide innovation.

\section{Protocol Specification}

\subsection{Data Types}

The protocol enforces a basic type system:

\begin{lstlisting}[language=,caption={Supported Data Types}]
string    - UTF-8 encoded text
uint      - Unsigned integer (0 to 2^256-1)
int       - Signed integer (-2^255 to 2^255-1)
bytes     - Raw byte array
bool      - Boolean (true/false)
address   - XRPL address (rBase58 encoded)
array<T>  - Ordered list of type T
object    - Nested structure
\end{lstlisting}

\subsection{Schema Definition}

Schemas are the foundation of XCS credentials, defining the structure and types of data that credentials can contain. Each schema acts as a template that issuers use when creating credentials, ensuring consistent data formats across the ecosystem.

\subsubsection{Schema Structure}

A schema consists of field definitions mapping names to types:

\begin{lstlisting}[language=json,caption={Basic Schema Structure}]
{
  "fields": {
    "name": "string",
    "age": "uint",
    "verified": "bool",
    "skills": "array<string>",
    "address": {
      "street": "string",
      "city": "string",
      "postalCode": "uint"
    }
  }
}
\end{lstlisting}

\subsubsection{Schema Creation Process}

Creating a schema involves three steps:

\begin{enumerate}
    \item \textbf{Define the schema structure} with required fields and types
    \item \textbf{Submit registration transaction} via payment with schema in memo
    \item \textbf{Protocol computes UID} after transaction confirmation
\end{enumerate}

\subsubsection{UID Computation}

Schema UIDs are computed by the protocol indexer after the transaction is confirmed on-chain. This ensures deterministic and unique identifiers:

\begin{lstlisting}[caption={UID Calculation (by Protocol)}]
schema_uid = SHA256(
  canonical_json(schema_definition) ||
  issuer_address ||
  ledger_index ||
  transaction_index
)
\end{lstlisting}

The UID calculation includes:
\begin{itemize}
    \item \textbf{Canonical JSON}: Ensures consistent ordering of fields
    \item \textbf{Issuer address}: Prevents UID collisions between issuers
    \item \textbf{Ledger index}: The confirmed ledger containing the transaction
    \item \textbf{Transaction index}: Position within the ledger
\end{itemize}

Since ledger and transaction indices are only known after confirmation, the UID must be computed by the protocol indexer, not the issuer. This ensures global uniqueness and prevents manipulation.

\subsubsection{Schema Versioning}

Schemas are immutable once registered. Versioning works through new schema creation:

\begin{lstlisting}[language=json,caption={Schema Evolution}]
// Version 1 - Original schema
{
  "fields": {
    "name": "string",
    "employeeId": "uint"
  }
}
// UID: abc123... (computed after registration)

// Version 2 - Extended schema
{
  "extends": ["abc123..."],  // Reference v1 by UID
  "fields": {
    "department": "string",  // Add new field
    "email": "string"
  }
}
// UID: def456... (new UID computed after registration)
\end{lstlisting}

Each version receives a new UID after its registration transaction confirms. Applications choose which versions to support, enabling natural migration without breaking existing credentials.

\subsection{Schema Inheritance}

Inheritance enables schema composition and reuse, reducing redundancy and promoting standardization.

\subsubsection{Multiple Inheritance}

Schemas can inherit from multiple parents, combining their fields:

\begin{lstlisting}[language=json,caption={Multiple Inheritance Example}]
// Parent 1: Basic Identity
{
  "uid": "identity_base",  // Computed by protocol
  "fields": {
    "name": "string",
    "dateOfBirth": "uint"
  }
}

// Parent 2: Contact Info
{
  "uid": "contact_base",  // Computed by protocol
  "fields": {
    "email": "string",
    "phone": "string"
  }
}

// Child: Employee combining both
{
  "extends": ["identity_base", "contact_base"],
  "fields": {
    "employeeId": "uint",
    "department": "string"
  }
}
// Result: name, dateOfBirth, email, phone, employeeId, department
\end{lstlisting}

\subsubsection{Selective Exclusion}

The \texttt{exclude} mechanism allows fine-grained control over inherited fields:

\begin{lstlisting}[language=json,caption={Selective Field Exclusion}]
{
  "extends": ["parent_uid_1", "parent_uid_2"],
  "exclude": {
    "parent_uid_1": ["sensitiveField"],
    "parent_uid_2": ["unnecessaryField"]
  },
  "fields": {
    "newField": "string"
  }
}
\end{lstlisting}

\subsubsection{Inheritance Rules}

The protocol enforces clear inheritance semantics during indexing:

\begin{enumerate}
    \item \textbf{Exclusion first}: Fields are excluded before checking conflicts
    \item \textbf{Conflict detection}: Same field name from multiple parents causes schema invalidity
    \item \textbf{Child override}: Child fields replace inherited fields with same name
    \item \textbf{Validation}: Protocol indexer marks schema invalid if conflicts exist after exclusion
\end{enumerate}

\begin{lstlisting}[language=json,caption={Conflict Resolution Example}]
// Invalid: Conflict without exclusion
{
  "extends": ["schema_a", "schema_b"],  // Both define "id" field
  "fields": { "status": "string" }
}
// Result: INVALID (marked by indexer) - field conflict

// Valid: Conflict resolved via exclusion
{
  "extends": ["schema_a", "schema_b"],
  "exclude": {
    "schema_b": ["id"]  // Exclude conflicting field
  },
  "fields": { "status": "string" }
}
// Result: Valid - uses schema_a's "id" field
\end{lstlisting}

This inheritance system enables powerful composition patterns while maintaining simplicity and preventing ambiguity. The protocol indexer validates inheritance rules during processing, ensuring only valid schemas are available for use.

\subsection{Protocol Operations}

The XCS protocol builds directly on top of the XRPL Credentials amendment, leveraging its native credential lifecycle operations while adding a schema layer through standard XRPL payment transactions. This design maintains compatibility with existing XRPL infrastructure while enabling extensibility.

\subsubsection{Protocol Composition}

XCS operates at two levels:

\textbf{Schema Layer}: Uses XRPL payment transactions with specially formatted memos to register and discover schemas. The memo field carries schema definitions, prefixed with \texttt{xcs:} to ensure namespace isolation.

\textbf{Credential Layer}: Directly utilizes the native Credentials amendment operations (CredentialCreate, CredentialAccept, CredentialDelete) without modification, referencing schemas by their UIDs.

\vspace{1em}
This layered approach means that XCS does not require changes to XRPL itself, it only adds a coordination layer on top.

\subsubsection{Memo Prefixing Strategy}

All XCS protocol operations use the \texttt{xcs:} prefix in memo fields:

\begin{lstlisting}[caption={XCS Memo Format}]
Memo:
  Type: "xcs:schema_register"
  Data: {schema_definition}
\end{lstlisting}

The prefix serves critical functions:
\begin{itemize}
    \item \textbf{Namespace isolation}: Prevents collision with other protocols using memos
    \item \textbf{Indexing efficiency}: Protocol indexers can quickly filter XCS transactions
    \item \textbf{Protocol versioning}: Future versions could use \texttt{xcs:v2:} prefixes
    \item \textbf{Operation types}: Different operations like \texttt{xcs:schema\_update} can be added
\end{itemize}

Only transactions with properly prefixed memos are processed by XCS indexers, ensuring clean protocol boundaries and preventing interference from unrelated XRPL activity.


\subsubsection{Operation Types}

\subsubsection{Operation Types}

The protocol supports four core operations:

\begin{enumerate}
    \item \textbf{Schema Registration}: Define new immutable credential structures (requires \texttt{xcs:} memo)
    \item \textbf{Credential Creation}: Issue credentials based on schemas (requires \texttt{xcs:} memo for protocol tracking)
    \item \textbf{Credential Acceptance}: Accept issued credentials (optional \texttt{xcs:} memo)
    \item \textbf{Credential Revocation}: Revoke existing credentials (optional \texttt{xcs:} memo)
\end{enumerate}

Schemas are immutable once registered. Any changes require creating a new schema, potentially inheriting from the previous version. This ensures perfect reproducibility and eliminates versioning ambiguity.

The \texttt{xcs:} memo prefix is mandatory for schema operations and credential creation to ensure proper protocol indexing. For acceptance and revocation, the memo is optional since the protocol can determine membership from existing credential state.

\subsubsection{Schema Registration}

Creates a new schema in the protocol registry:

\begin{lstlisting}[caption={Schema Registration Transaction}]
Transaction Type: Payment
Destination: SCHEMA_REGISTRY_ADDRESS
Amount: SCHEMA_REGISTRATION_FEE XRP
Memo:
  Type: "xcs:schema_register"
  Data: {
    "extends": [optional_parent_uids],
    "exclude": {optional_exclusions},
    "fields": {field_definitions}
  }
\end{lstlisting}

\subsubsection{Credential Issuance}

Issues a credential conforming to a registered schema:

\begin{lstlisting}[caption={Credential Creation}]
Transaction Type: CredentialCreate
Fields:
  Subject: recipient_address
  CredentialType: schema_uid
  URI: storage_location
  Expiration: optional_unix_timestamp
Memo:
  Type: "xcs:credential_create"
\end{lstlisting}


\subsubsection{Credential Acceptance}

Accepts an issued credential:

\begin{lstlisting}[caption={Credential Acceptance}]
Transaction Type: CredentialAccept
Fields:
  Issuer: issuer_address
  CredentialType: schema_uid
Memo (Optional):
  Type: "xcs:credential_accept"
  Data: {
    "acceptance_context": "optional_metadata"
  }
\end{lstlisting}

\subsubsection{Credential Revocation}

Revokes an existing credential:

\begin{lstlisting}[caption={Credential Deletion}]
Transaction Type: CredentialDelete
Fields:
  Subject: subject_address
  Issuer: issuer_address
  CredentialType: schema_uid
Memo (Optional):
  Type: "xcs:credential_revoke"
  Data: {
    "reason": "optional_revocation_reason",
    "timestamp": unix_timestamp
  }
\end{lstlisting}

While the memo is optional, including revocation reasons can help with audit trails and compliance requirements.

\subsubsection{Future Operations}

The protocol's extensible design allows for future operations without breaking changes for example:

\begin{itemize}
    \item \texttt{xcs:schema\_deprecate}: Mark schemas as deprecated
    \item \texttt{xcs:credential\_suspend}: Temporary credential suspension
    \item \texttt{xcs:batch\_operation}: Bulk operations for efficiency
    \item \texttt{xcs:delegation\_grant}: Delegate issuance authority
\end{itemize}

Each new operation would use the \texttt{xcs:} prefix, ensuring backward compatibility and clean protocol evolution.
\subsection{Storage Patterns}

While storage is not part of the protocol, common patterns emerge:

\textbf{Public Credentials:}
\begin{center}
XRPL Entry $\rightarrow$ IPFS Content $\rightarrow$ Public Verification\\
URI: \texttt{ipfs://QmXxx...}
\end{center}

\textbf{Private Credentials:}
\begin{center}
XRPL Entry $\rightarrow$ Private API $\rightarrow$ Authorized Verification\\
URI: \texttt{https://issuer.com/api/credential/\{id\}}
\end{center}

\subsection{Verification Flow}

\begin{enumerate}
    \item \textbf{Retrieve credential from XRPL}
    \begin{itemize}
        \item Query by subject, issuer, and credential type
        \item Check acceptance status and expiration
    \end{itemize}
    
    \item \textbf{Fetch credential data}
    \begin{itemize}
        \item Public: Retrieve from IPFS/public storage
        \item Private: Query issuer API with authorization
    \end{itemize}
    
    \item \textbf{Validate against schema}
    \begin{itemize}
        \item Retrieve schema from registry
        \item Recursively resolve parent schemas
        \item Type-check credential data
    \end{itemize}
    
    \item \textbf{Application-specific verification}
    \begin{itemize}
        \item Additional business logic
        \item Privacy-preserving proofs
        \item External data correlation
    \end{itemize}
\end{enumerate}

\section{Implementation Considerations}

\subsection{Schema Registry}

The schema registry (\texttt{SCHEMA\_REGISTRY\_ADDRESS}) can be:
\begin{itemize}
    \item A burn address for truly decentralized operation
    \item A managed address that funds infrastructure
    \item A multi-signature address for community governance
\end{itemize}

The registry maintains an index of all schemas, enabling:
\begin{itemize}
    \item Discovery by UID
    \item Browsing by issuer
    \item Search by inheritance tree
\end{itemize}

\subsection{Type Validation}

Type validation occurs at multiple layers:
\begin{enumerate}
    \item \textbf{Client-side}: Immediate feedback during credential creation
    \item \textbf{Indexer}: Validation during registry indexing
    \item \textbf{Verifier}: Final validation during credential verification
\end{enumerate}

Invalid schemas are marked but not removed, maintaining immutability.

\subsection{Batch Operations}

The upcoming XRPL batch transactions feature enables:
\begin{itemize}
    \item Multiple credential issuance in one transaction
    \item Atomic schema registration with initial credentials
    \item Bulk revocation for compromised issuers
\end{itemize}

No special protocol support needed; native XRPL features suffice.

\subsection{Performance Considerations}

\begin{itemize}
    \item \textbf{Schema caching}: Frequently used schemas should be cached
    \item \textbf{Inheritance resolution}: Build materialized views of inherited fields
    \item \textbf{IPFS pinning}: Critical credentials need reliable pinning services
    \item \textbf{API rate limiting}: Private credential APIs need abuse protection
\end{itemize}

\section{Use Cases}

\subsection{KYC/AML Compliance}

Financial institutions can:
\begin{enumerate}
    \item Define schemas for required compliance data
    \item Accept credentials from trusted KYC providers
    \item Verify credentials without accessing personal data
    \item Maintain audit trails on-chain
\end{enumerate}

\subsection{Gaming Achievements}

Games can issue credentials for:
\begin{itemize}
    \item Player achievements and levels
    \item Item ownership and rarity
    \item Tournament participation
    \item Cross-game reputation
\end{itemize}

\subsection{Academic Credentials}

Educational institutions can:
\begin{itemize}
    \item Issue tamper-proof diplomas and certificates
    \item Enable instant verification by employers
    \item Support micro-credentials and continuous learning
    \item Facilitate credit transfer between institutions
\end{itemize}

\subsection{XRP Name Service}

A decentralized naming system using XCS:

\begin{lstlisting}[language=json,caption={XRP Name Service Schema}]
{
  "schema": {
    "fields": {
      "name": "string",
      "address": "address",
      "metadata": {
        "avatar": "string",
        "bio": "string",
        "social": "array<string>"
      },
      "subdomains": "array<object>"
    }
  }
}
\end{lstlisting}

Benefits:
\begin{itemize}
    \item Human-readable addresses
    \item Decentralized profile data
    \item Hierarchical namespace support
    \item Integration with existing XRPL features
\end{itemize}

\section{Comparison with Existing Systems}

\subsection{XCS vs. W3C Verifiable Credentials}

\begin{table}[h]
\centering
\begin{tabular}{lll}
\toprule
\textbf{Aspect} & \textbf{XCS} & \textbf{W3C VC} \\
\midrule
Philosophy & Minimal protocol, rich applications & Comprehensive standard \\
Storage & Flexible (IPFS, APIs, etc.) & Typically includes data model \\
Privacy & Application layer & Built-in (JSON-LD, ZKP) \\
Complexity & Low & High \\
Interoperability & Via simple types & Via JSON-LD contexts \\
\bottomrule
\end{tabular}
\caption{Comparison with W3C Verifiable Credentials}
\end{table}

\subsection{XCS vs. Other Blockchain Credentials}

\begin{table}[h]
\centering
\begin{tabular}{p{3cm}p{5cm}p{5cm}}
\toprule
\textbf{Platform} & \textbf{Approach} & \textbf{Trade-offs} \\
\midrule
Ethereum/ERC-721 & NFT-based & High gas costs, complex smart contracts \\
Hyperledger Indy & Purpose-built & Separate blockchain required \\
Bitcoin Ordinals & Inscriptions & Limited programmability \\
XCS & Native XRPL & Low fees, simple model, composable \\
\bottomrule
\end{tabular}
\caption{Comparison with Other Blockchain Credential Systems}
\end{table}

\section{Security Considerations}

\subsection{Schema Poisoning}

Malicious actors might create schemas that appear legitimate. 

\textbf{Mitigation:}
\begin{itemize}
    \item Schema explorer shows issuer reputation
    \item Applications maintain trusted issuer lists
    \item Community-curated schema registries
\end{itemize}

\subsection{Storage Availability}

IPFS content may become unavailable. 

\textbf{Mitigation:}
\begin{itemize}
    \item Multiple pinning services
    \item Fallback to issuer-hosted data
    \item Critical data redundancy requirements
\end{itemize}

\subsection{Privacy Leakage}

On-chain credential existence reveals relationships. 

\textbf{Mitigation:}
\begin{itemize}
    \item Use separate accounts for different contexts
    \item Batch credentials to obscure timing
    \item Private credentials for sensitive relationships
\end{itemize}

\subsection{Type Confusion}

Incorrect type validation could enable attacks. 

\textbf{Mitigation:}
\begin{itemize}
    \item Strict type checking at all layers
    \item Canonical JSON representation
    \item Schema versioning for upgrades
\end{itemize}

\section{Future Work}

\subsection{Protocol Extensions}

While keeping the core protocol minimal, future extensions could include:

\begin{itemize}
    \item \textbf{Delegation}: Allow issuers to delegate issuance authority
    \item \textbf{Conditional Credentials}: Credentials that activate based on conditions
    \item \textbf{Cross-chain Bridges}: Interoperability with other blockchain credentials
\end{itemize}

\subsection{Ecosystem Development}

\begin{itemize}
    \item \textbf{SDK Development}: Libraries for major programming languages
    \item \textbf{Schema Marketplace}: Curated, reusable schema templates
    \item \textbf{Verification Services}: Third-party verification aggregators
    \item \textbf{Privacy Layers}: ZK-SNARK/STARK integration libraries
\end{itemize}

\subsection{Governance}

As the protocol matures:
\begin{itemize}
    \item Establish a standards body for schema conventions
    \item Create certification programs for issuers
    \item Develop dispute resolution mechanisms
\end{itemize}

\section{Conclusion}

XCS demonstrates that effective protocols embrace constraints. By deliberately limiting scope to essential coordination primitives—schema registration, credential lifecycle, and type systems—XCS achieves what comprehensive standards cannot: simplicity that enables innovation.

The protocol's minimalism is not a limitation but its greatest strength. Developers can build KYC systems, gaming achievements, academic credentials, or entirely novel applications without protocol changes. The thin protocol/rich application model ensures XCS can evolve with use cases rather than trying to predict them.

As the credential ecosystem grows increasingly complex, XCS provides a stable foundation: simple enough to implement, flexible enough to extend, and robust enough to build upon. The protocol does not solve every problem—it provides the tools for others to solve their problems.

\section*{References}

\begin{enumerate}[label={[\arabic*]}]
    \item XRPL Credentials Amendment Documentation
    \item W3C Verifiable Credentials Data Model
    \item Decentralized Identifiers (DIDs) Specification
    \item IPFS Protocol Specification
    \item XRP Ledger Protocol Documentation
\end{enumerate}

\appendix

\section{Example Implementations}

\subsection{Simple KYC Credential}

Schema:
\begin{lstlisting}[language=json]
{
  "fields": {
    "verified": "bool",
    "level": "uint",
    "country": "string",
    "expiry": "uint"
  }
}
\end{lstlisting}

\subsection{Composable Identity}

Base Identity:
\begin{lstlisting}[language=json]
{
  "uid": "base_identity_123",
  "fields": {
    "name": "string",
    "dateOfBirth": "uint"
  }
}
\end{lstlisting}

Extended KYC:
\begin{lstlisting}[language=json]
{
  "extends": ["base_identity_123"],
  "fields": {
    "nationality": "string",
    "taxId": "string"
  }
}
\end{lstlisting}

\subsection{XRP Name Service Entry}

\begin{lstlisting}[language=json]
{
  "fields": {
    "name": "alice.xrp",
    "address": "rN7n7otQDd6FczFgLdAtqCSMjqXqLSRgDc",
    "metadata": {
      "avatar": "ipfs://Qm...",
      "bio": "Blockchain developer",
      "social": ["twitter:@alice", "github:alice"]
    },
    "subdomains": [
      {
        "name": "pay",
        "address": "rDifferentAddressForPayments..."
      }
    ]
  }
}
\end{lstlisting}

\vspace{2cm}
\begin{center}
\rule{0.5\textwidth}{0.4pt}\\[0.5cm]
\textit{XCS Protocol v1.0 - Draft Specification}
\end{center}

\end{document}
